% Patch 0615 A5_F: Second-Order Closed-Form Substrate-Current Coefficients
% at Face-Aligned Reading C in the 600-Cell
%
% OPEN-FP-F1-5 Sequence-2B closure artifact (THEO-DSL-9 candidate). Computes the
% O(delta^2) face-aligned coefficients in the residual-symmetry-invariant subspace
% at the host vertex, using the same Mechanism-A 2-edge-path assembly as the vertex
% (THEO-DSL-5) and edge (THEO-DSL-7) closures, with n = n_Fperp.
%
% This artifact also registers an ANTI-ERASURE CORRECTION to THEO-DSL-8 (Patch 0597):
% the true residual symmetry of (v_host + n_Fperp) is the Klein four-group V_4 of
% ORDER 4, not C_s of order 2; the invariant subspace is the 2D span{n_rho, n_ax},
% not the 3D (u_i - u_j)^perp. Consequently the dimensional-growth pattern across
% the three Reading-C variants is 1 -> 2 -> 2 (I_h -> D_5 -> V_4), not 1 -> 2 -> 3.
%
% Sketch-document Layer 3 under (H5_F) (same path-class weight ansatz W(P)=1 as the
% vertex THEO-DSL-5 and edge THEO-DSL-7), per DG-F15-4 default-(A) per-variant scope
% qualifier and DG-F15-5 / DG-F15A-5 default-(A) single-theorem registration.
%
% CHANGELOG
%   v1.0 (Patch 0615, Session 147, 28 May 2026): closure of OPEN-FP-F1-5 Sequence-2B.
%     Reproduces the vertex alpha_2 = -9/phi^2 cross-check, verifies the V_4 residual
%     symmetry directly on the 600-cell vertex set, and computes
%       alpha_2^(rho) = -14 + 7 phi = -7/phi^2 ~ -2.674,
%       alpha_2^(ax)  =  -6 + 2 phi           ~ -2.764,
%     with the two non-invariant components (n_diff via sigma_E; n_perp3 via
%     sigma_face) vanishing identically. Verify script:
%     code/verify_face_alpha2_closure.py. Reasoning: hardened_theorems/reasoning/0615.md.
% Version 1.0 --- 17 August 2026 (Patch 3212):
%   added the generated plain-language appendix glossing the internal
%   programme identifiers used in this paper (founder jargon ruling, 16
%   Aug 2026). No claim, equation, number or section changed.
%   This is the FIRST version stamp assigned to this paper; 1.0 denotes
%   the first stamped revision, NOT a claim of v1.0 shipped grade.

\documentclass[11pt]{article}
\usepackage[utf8]{inputenc}
\usepackage{amsmath,amssymb,amsthm}
\usepackage[margin=1in]{geometry}
\usepackage{hyperref}
\usepackage{booktabs}

\theoremstyle{plain}
\newtheorem{theorem}{Theorem}
\newtheorem{lemma}{Lemma}
\newtheorem{corollary}{Corollary}
\theoremstyle{definition}
\newtheorem*{remark}{Remark}
\newtheorem*{antierasure}{Anti-erasure Remark}

\newcommand{\nhat}{\hat{n}}
\newcommand{\nrho}{\hat{n}_\rho}
\newcommand{\nax}{\hat{n}_{\text{ax}}}
\newcommand{\nperpF}{\hat{n}_{F\perp}}
\newcommand{\nperpiii}{\hat{n}_{\perp 3}}
\newcommand{\ndiff}{\hat{n}_{\text{diff}}}
\newcommand{\vhost}{v_{\text{host}}}
\newcommand{\ui}{u_i}
\newcommand{\uj}{u_j}
\newcommand{\jvec}{\vec{j}}
\newcommand{\Reals}{\mathbb{R}}
\newcommand{\Icoh}{I_h}
\newcommand{\Vfour}{V_4}
\newcommand{\Cs}{C_s}
\newcommand{\Dfive}{D_5}

\title{Second-Order Closed-Form Substrate-Current Coefficients\\
at Face-Aligned Reading C in the 600-Cell\\
\large{(OPEN-FP-F1-5 Sequence-2B Closure Artifact A5$_F$; THEO-DSL-9 candidate,\\
$\alpha_2^{(\rho)} = -7/\phi^2$, $\alpha_2^{(\text{ax})} = 2\phi - 6$;
with a $\Cs \to \Vfour$ residual-symmetry correction to THEO-DSL-8)}}
\author{Thomas Lee Abshier, ND, and Claude Opus 4.7\\ Hyperphysics Institute --- F.1 Dynamical Substrate Law trajectory}
\date{28 May 2026 (Session 147, CPP Patch 0615)\\[2pt]
{\small Version~1.0 --- 17 August 2026 (Patch 3212: added the generated plain-language appendix glossing the internal programme identifiers used in this paper (founder jargon ruling, 16 Aug 2026). No claim, equation, number or section changed.)}}

\begin{document}
\maketitle

\begin{abstract}
We close the OPEN-FP-F1-5 Sequence-2B trajectory: the second-order ($\mathcal{O}(\delta^2)$)
closed-form coefficients of the net DI-bit substrate current at the host vertex of the
600-cell under \emph{face-aligned} Reading C, computed with the identical Mechanism-A
two-edge-path assembly used for the vertex-aligned (THEO-DSL-5, Patch 0591) and
edge-aligned (THEO-DSL-7, Patch 0613) closures, specialised to $\nhat = \nperpF$, the
face-centroid (substrate-primitive) direction. The computation yields, exact in
$\mathbb{Q}[\phi]$,
\[
\jvec_2^{\text{face}}(\vhost) \;=\; \alpha_2^{(\rho)}\,\nrho \;+\; \alpha_2^{(\text{ax})}\,\nax,
\qquad
\alpha_2^{(\rho)} = -14 + 7\phi = -\frac{7}{\phi^2} \approx -2.674,
\quad
\alpha_2^{(\text{ax})} = -6 + 2\phi \approx -2.764,
\]
with \emph{both} components negative (like vertex-aligned, unlike the positive
edge-aligned coefficients). In the course of the closure we discovered, and verify
directly on the 600-cell vertex set, that the residual symmetry of
$(\vhost + \nperpF)$ is the Klein four-group $\Vfour$ of \textbf{order 4}, not the
$\Cs$ of order 2 recorded in THEO-DSL-8 (Patch 0597): all three non-trivial $\Vfour$
elements fix the centroid direction $\nperpF$ and hence preserve Mechanism A's rate
function. The $\Vfour$-invariant subspace is therefore the \textbf{2-dimensional}
$\operatorname{span}\{\nrho, \nax\}$, not the 3-dimensional $(\ui - \uj)^\perp$;
independently, THEO-DSL-8's claimed three-vector basis $\{\nrho, \nax, \nperpF\}$ is
degenerate ($\nperpF \in \operatorname{span}\{\nrho, \nax\}$). The corrected
dimensional-growth pattern across the three Reading-C variants is therefore
$1 \to 2 \to 2$ (as the stabiliser runs $\Icoh \to \Dfive \to \Vfour$), not the
previously-recorded $1 \to 2 \to 3$. We register this as an anti-erasure correction
(\S\ref{sec:antierasure}) and recommend a follow-on structural revision of the
THEO-DSL-8 artifact under multi-reviewer scrutiny. All claims are verified
numerically and symbolically over $\mathbb{Q}[\phi]$ by
\texttt{code/verify\_face\_alpha2\_closure.py}.
\end{abstract}

\section{Setup}\label{sec:setup}

We preserve the vertex-host convention of the OPEN-FP-F1-5 trajectory (Patches 0596,
0597): $\vhost$ is a 600-cell vertex; only the direction of the substrate primitive
$\nhat$ varies among the three Reading-C variants. Vectors are taken at unit
circumradius ($|v|=1$ for every vertex; edge length $1/\phi$, $\phi = (1+\sqrt5)/2$).

\paragraph{Face and primitive direction.} Choose a triangular face
$F = \{\vhost, \ui, \uj\}$ incident to $\vhost$, where $\ui, \uj$ are first-shell
neighbours of $\vhost$ that are themselves adjacent (so $\vhost\cdot\ui = \vhost\cdot\uj
= \ui\cdot\uj = \phi/2$). There are $30$ such faces incident to each vertex. The
face-aligned substrate primitive is the centroid direction
\begin{equation}\label{eq:nperp}
\nperpF \;=\; \frac{\vhost + \ui + \uj}{\lVert \vhost + \ui + \uj\rVert},
\end{equation}
which is perpendicular to the face plane and is the rate-function direction of
(H5$_F$): $r(\hat e) = r_0(1 + \delta\,\hat e\cdot\nperpF)$.

\paragraph{Frame at $\vhost$.} We use the orthonormal frame
\begin{align}
\nrho &= \vhost && \text{(radial)},\\
\nax  &= \frac{P_{\perp\vhost}(\ui + \uj - 2\vhost)}{\lVert\cdot\rVert}
         = \frac{P_{\perp\vhost}(\ui + \uj)}{\lVert\cdot\rVert}
         && \text{(in-face axial, tangent to $S^3$)},\\
\nperpiii &= \text{unit normal to } \operatorname{span}\{\vhost, \ui, \uj\}
         && \text{(out-of-3-flat)},\\
\ndiff &= \frac{\ui - \uj}{\lVert\ui - \uj\rVert} && \text{(within-face edge)},
\end{align}
where $P_{\perp\vhost}(x) = x - (x\cdot\vhost)\vhost$. These four directions are
mutually orthonormal and span $\Reals^4$ at $\vhost$. The centroid direction of
\eqref{eq:nperp} is \emph{not} independent of this frame:
\begin{equation}\label{eq:nFdegenerate}
\ui + \uj = \phi\,\vhost + P_{\perp\vhost}(\ui+\uj)
\;\Longrightarrow\;
\nperpF \;\propto\; (1+\phi)\,\nrho + \lVert P_{\perp\vhost}(\ui+\uj)\rVert\,\nax
\;\in\; \operatorname{span}\{\nrho, \nax\}.
\end{equation}

\section{Residual symmetry is $\Vfour$, not $\Cs$}\label{sec:v4}

\begin{lemma}[$\Vfour$ residual symmetry]\label{lem:v4}
The stabiliser in $\Icoh = \operatorname{stab}_{H_4}(\vhost)$ of the pair
$(\vhost,\,\nperpF)$ --- equivalently of the host vertex together with the centroid
direction of the face $F$ --- is the Klein four-group
\[
\Vfour \;=\; \{e,\ \sigma_E,\ \sigma_{\!F},\ C_2\}, \qquad |\Vfour| = 4,
\]
where $\sigma_E$ is the reflection through the hyperplane $(\ui-\uj)^\perp$ (swaps
$\ui\leftrightarrow\uj$, flips $\ndiff$), $\sigma_{\!F}$ is the reflection through the
hyperplane $\operatorname{span}\{\vhost,\ui,\uj\}$ (fixes $\vhost,\ui,\uj$, flips
$\nperpiii$), and $C_2 = \sigma_E\sigma_{\!F}$. All three non-trivial elements fix
$\nperpF$.
\end{lemma}

\begin{proof}
The unordered-face stabiliser in $\Icoh$ has order $|\Icoh|/30 = 120/30 = 4$ by
orbit--stabiliser, and is $\Vfour$ (Klein four), as recorded in THEO-DSL-8 \S2.2.
Each of $e, \sigma_E, \sigma_{\!F}, C_2$ fixes $\vhost$ and permutes $\{\ui,\uj\}$
among themselves, hence fixes the centroid $\vhost+\ui+\uj$ and therefore the
direction $\nperpF$ of \eqref{eq:nperp}. In particular $\sigma_{\!F}$ fixes $\ui$ and
$\uj$ individually, so it fixes their sum and $\nperpF$. Because $\nperpF$ is fixed by
all of $\Vfour$, the rate function $r(\hat e) = r_0(1+\delta\,\hat e\cdot\nperpF)$ is
$\Vfour$-invariant. (Direct verification: $\sigma_E, \sigma_{\!F}, C_2$, realised as
the orthogonal matrices $I - 2\ndiff\ndiff^\top$, $I - 2\nperpiii\nperpiii^\top$, and
their product, each map the 120-vertex set to itself and satisfy $R\vhost=\vhost$,
$R\nperpF=\nperpF$; see the verify script.)
\end{proof}

\begin{lemma}[Two-dimensional invariant subspace]\label{lem:2d}
The $\Vfour$-invariant subspace of $\Reals^4$ at $\vhost$ is the $2$-dimensional
$\operatorname{span}\{\nrho, \nax\} = (\ui-\uj)^\perp \cap \operatorname{span}\{\vhost,\ui,\uj\}$.
\end{lemma}

\begin{proof}
$\sigma_E$ acts as $+1$ on $(\ui-\uj)^\perp$ (dimension 3) and $-1$ on $\Reals\ndiff$;
$\sigma_{\!F}$ acts as $+1$ on $\operatorname{span}\{\vhost,\ui,\uj\}$ (dimension 3)
and $-1$ on $\Reals\nperpiii$. A vector fixed by both lies in
$(\ui-\uj)^\perp \cap \operatorname{span}\{\vhost,\ui,\uj\}$. The latter equals
$\operatorname{span}\{\vhost,\ \ui+\uj\} = \operatorname{span}\{\nrho,\nax\}$,
which is $2$-dimensional (the Gram determinant of $\{\vhost,\ui,\uj\}$ is
$(1-\phi/2)^2(1+\phi)>0$, so the three vectors are independent and their span is a
genuine $3$-flat with $1$-dimensional normal $\nperpiii$). A direct averaging
projector over the four $\Vfour$ elements has rank $2$ (verify script).
\end{proof}

\begin{lemma}[$\Vfour$-equivariance of the current, all orders]\label{lem:equiv}
For every $g\in\Vfour$ and every $k\ge 0$,
$g\cdot\jvec_k^{\text{face}}(\vhost) = \jvec_k^{\text{face}}(\vhost)$, hence
$\jvec_k^{\text{face}}(\vhost) \in \operatorname{span}\{\nrho,\nax\}$ for all $k\ge1$.
\end{lemma}

\begin{proof}
Identical in structure to THEO-DSL-6 Lemma~1 / THEO-DSL-8 Lemma~1, now with the full
$\Vfour$ in place of $\Cs$: (i) $\Vfour$-equivariance of the rate function
(Lemma~\ref{lem:v4}, $g\nperpF=\nperpF$); (ii) $\Vfour$ permutes the path set
$\mathcal{P}_k(\vhost)$ as graph automorphisms fixing $\vhost$ (shell-locality (H4));
(iii) $O(4)$-equivariance of edge-direction vectors. Re-indexing the order-$k$ sum by
the $g$-action gives $g\cdot\jvec_k^{\text{face}} = \jvec_k^{\text{face}}$, so
$\jvec_k^{\text{face}}$ lies in the $\Vfour$-invariant subspace, which is
$\operatorname{span}\{\nrho,\nax\}$ by Lemma~\ref{lem:2d}.
\end{proof}

\section{Second-order coefficients}\label{sec:coeff}

The $\mathcal{O}(\delta^2)$ coefficient is the path-integral assembly over the $144$
directed two-edge paths $\vhost \to \ui \to v'$:
\begin{equation}\label{eq:assembly}
\jvec_2^{\text{face}}(\vhost)
\;=\; \sum_{P=(\vhost\to u\to v')} (\hat e_1\cdot\nperpF)(\hat e_2\cdot\nperpF)\,\hat e_1,
\qquad \hat e_1 = \widehat{u-\vhost},\ \ \hat e_2 = \widehat{v'-u},
\end{equation}
with path-class weight $W(P)=1$ per (H5$_F$) at sketch-document Layer 3 (the same
ansatz as THEO-DSL-5/-7). Carrying out \eqref{eq:assembly} and projecting onto the
frame gives the closed forms

\begin{theorem}[Face-aligned $\mathcal{O}(\delta^2)$ closure]\label{thm:main}
Under (H1$_F$)--(H5$_F$) (the THEO-DSL-8 hypothesis stack with (H5$_F$) the
$W(P)=1$ path-class weight ansatz), the second-order face-aligned substrate-current
coefficient is, exact in $\mathbb{Q}[\phi]$,
\[
\jvec_2^{\text{face}}(\vhost)
= \Bigl(-14 + 7\phi\Bigr)\nrho + \Bigl(-6 + 2\phi\Bigr)\nax,
\qquad
\alpha_2^{(\rho)} = -14+7\phi = -\tfrac{7}{\phi^2}\approx-2.674,
\quad
\alpha_2^{(\text{ax})} = -6+2\phi \approx -2.764,
\]
with $\alpha_2^{(\perp 3)} = 0$ (forced by $\sigma_{\!F}$) and
$\alpha_2^{(\text{diff})} = 0$ (forced by $\sigma_E$; the THEO-DSL-8 conclusion).
The result is independent of the choice of incident face (verified over all $30$).
\end{theorem}

The four path classes by endpoint shell decompose as in Table~\ref{tab:classes}
(basis $\{\nrho,\nax,\nperpiii\}$; the $\nperpiii$ column is identically zero,
class by class).

\begin{table}[h]\centering
\begin{tabular}{lccc}
\toprule
class & \# paths & coeff $\nrho$ & coeff $\nax$ \\
\midrule
B.1 (round-trip, $v'=\vhost$) & 12 & $-\tfrac{7}{6} + \phi$ & $-\tfrac13 + \tfrac23\phi$ \\
B.2 (first-shell, $v'\in S_1$) & 60 & $-5 + \tfrac{10}{3}\phi$ & $-\tfrac53 + \tfrac53\phi$ \\
B.3 (cross-shell, $v'\in S_2$) & 60 & $-\tfrac{20}{3} + \tfrac52\phi$ & $-\tfrac{10}{3}$ \\
B.4 (third-shell, $v'\in S_3$) & 12 & $-\tfrac{7}{6} + \tfrac16\phi$ & $-\tfrac23 - \tfrac13\phi$ \\
\midrule
\textbf{TOTAL} & 144 & $-14 + 7\phi$ & $-6 + 2\phi$ \\
\bottomrule
\end{tabular}
\caption{Per-shell-class decomposition of $\jvec_2^{\text{face}}$. Every class has
identically zero $\nperpiii$ component (and zero $\ndiff$ component), consistent with
Lemma~\ref{lem:equiv}.}
\label{tab:classes}
\end{table}

\paragraph{Cross-variant comparison.} The three Reading-C $\mathcal{O}(\delta^2)$
closures now read: vertex $\alpha_2 = -9/\phi^2 \approx -3.438$ (1D, THEO-DSL-5);
edge $(\alpha_2^{(\rho)},\alpha_2^{(\text{edge})}) = (9/(2\phi),\,9\phi-12) \approx
(+2.781,\,+2.562)$ (2D, THEO-DSL-7); face
$(\alpha_2^{(\rho)},\alpha_2^{(\text{ax})}) = (-7/\phi^2,\,2\phi-6) \approx
(-2.674,\,-2.764)$ (2D, this artifact). The radial coefficient $-7/\phi^2$ at
face-aligned is structurally close to the vertex $-9/\phi^2$ (same sign, same
$1/\phi^2$ scale), while the edge case sits apart in sign --- a consequence of the
non-vanishing first-shell-transverse (B.2) class structure that differs across the
three stabilisers.

\section{Mandatory vertex cross-check}\label{sec:crosscheck}

Running the same assembly \eqref{eq:assembly} with $\nhat = \vhost$ (vertex-aligned)
reproduces $\alpha_2 = -9/\phi^2$ exactly, with perpendicular residual
$\sim 10^{-16}$ (verify script \S2). This guarantees the machinery is calibrated; the
face-aligned coefficients are computed by the identical code path with only
$\nhat\to\nperpF$ and the decomposition basis changed.

\section{Anti-erasure correction to THEO-DSL-8}\label{sec:antierasure}

\begin{antierasure}
THEO-DSL-8 (Patch 0597) records, for the face-aligned variant, the residual symmetry
$\Cs$ of order 2 and a $3$-dimensional invariant subspace $(\ui-\uj)^\perp =
\operatorname{span}\{\nrho,\nax,\nperpF\}$. Both are incorrect, by two independent
errors, surfaced here per the programme's anti-erasure discipline (cf.\ the
$D_3\to\Dfive$ and $D_2\to\Cs$ stabiliser corrections at Patches 0596--0597):
\begin{enumerate}
\item \textbf{Residual symmetry $\Cs \to \Vfour$ (order $2 \to 4$).} THEO-DSL-8 \S2.2
(``Subgroup preserving oriented $\nperpF$'') asserts that two of the four $\Vfour$
elements flip $\nperpF\to-\nperpF$. This is false: $\nperpF$ is the face \emph{centroid}
direction \eqref{eq:nperp}, fixed by every element of the face stabiliser (each
permutes $\{\ui,\uj\}$ and fixes $\vhost$). In particular the reflection
$\sigma_{\!F}$ through $\operatorname{span}\{\vhost,\ui,\uj\}$ fixes $\ui,\uj,\vhost$
individually and hence fixes $\nperpF$; it is a genuine 600-cell automorphism
(verified). The correct stabiliser of $(\vhost,\nperpF)$ is $\Vfour$ of order 4
(Lemma~\ref{lem:v4}).
\item \textbf{Invariant subspace $3$D $\to$ $2$D, and a degenerate basis.} With the
correct $\Vfour$, the invariant subspace is $2$-dimensional
$\operatorname{span}\{\nrho,\nax\}$ (Lemma~\ref{lem:2d}): $\sigma_E$ removes $\ndiff$
and $\sigma_{\!F}$ removes $\nperpiii$. Independently, THEO-DSL-8's claimed basis
$\{\nrho,\nax,\nperpF\}$ for its (over-stated) $3$D subspace is linearly
\emph{dependent}: by \eqref{eq:nFdegenerate}, $\nperpF\in\operatorname{span}\{\nrho,\nax\}$,
so those three vectors span only $2$D. ``Perpendicular to the face plane'' is a
$2$-dimensional condition; the specific centroid choice $\nperpF$ happens to lie in
$\operatorname{span}\{\nrho,\nax\}$, and the genuinely independent face-normal is
$\nperpiii$ (the out-of-3-flat direction), which is \emph{excluded} by $\sigma_{\!F}$.
\end{enumerate}
\textbf{Net effect.} The face-aligned structural conclusion strengthens (smaller
invariant subspace, $2$D not $3$D), and the dimensional-growth pattern across the
Reading-C variants is $1 \to 2 \to 2$ (stabilisers $\Icoh \to \Dfive \to \Vfour$),
\emph{not} $1 \to 2 \to 3$. The THEO-DSL-8 \emph{theorem statement}
($\jvec_k\perp(\ui-\uj)$) remains \emph{true} but is non-tight; the stronger true
statement is $\jvec_k\in\operatorname{span}\{\nrho,\nax\}$. We do not silently rewrite
the THEO-DSL-8 entry here; we register the correction and recommend a follow-on
structural-revision Patch of \texttt{face\_aligned\_invariant\_subspace\_structural.tex}
under multi-reviewer scrutiny, propagating the $1\to2\to2$ pattern to all programme
records that currently carry $1\to2\to3$.
\end{antierasure}

\section{Hypotheses, exclusions, falsifier}\label{sec:scope}

\paragraph{Hypotheses.} (H1$_F$) face-aligned Reading C, $\nhat=\nperpF$ of
\eqref{eq:nperp}; (H2) vertex-host unit-circumradius 600-cell; (H3$_F$) residual
symmetry $\Vfour$ at $\vhost$ (corrected; Lemma~\ref{lem:v4}); (H4) F.1 Theorem 6.1
perturbation-locality (THEO-DSL-1); (H5$_F$) Mechanism-A path-integral extension at
$\mathcal{O}(\delta^k)$ with path-class weight $W(P)=1$, sketch-document Layer 3
(DG-F15-4 default-(A)), the load-bearing scope qualifier.

\paragraph{Exclusion classes.}
(E1) $\mathcal{O}(\delta^k)$ for $k\ge 3$ not computed.
(E2) Vertex-aligned (THEO-DSL-5) and edge-aligned (THEO-DSL-7) variants are separate.
(E3) (H5$_F$) at sketch-document Layer 3 per DG-1 sub-option (A); publication-grade
unconditional promotion deferred to priority queue item (G).
(E4) Layer-4 derivation of the 600-cell substrate (OPEN-FP-F1-2) not in scope.
(E5) Whether $\alpha_2^{(\rho)}$ and $\alpha_2^{(\text{ax})}$ are further constrained
at framework-axiom level (beyond the $\Vfour$ admission) is open.

\paragraph{Falsifier.} Any first-principles $\Vfour$-equivariant recomputation under
(H1$_F$)--(H5$_F$) yielding $\jvec_2^{\text{face}}\notin\operatorname{span}\{\nrho,\nax\}$,
or yielding $(\alpha_2^{(\rho)},\alpha_2^{(\text{ax})})\neq(-7/\phi^2,\,2\phi-6)$, or a
vertex cross-check not reproducing $-9/\phi^2$; or a demonstration that $\sigma_{\!F}$
is not a 600-cell automorphism fixing $\nperpF$ (which would restore $\Cs$/3D).

\paragraph{Layer rigour.} Sketch-document Layer 3 under (H5$_F$), the same status as
the vertex THEO-DSL-5 and edge THEO-DSL-7 coefficient closures. The $\Vfour$
residual-symmetry and $2$D invariant-subspace results (Lemmas~\ref{lem:v4}--\ref{lem:2d})
are publication-grade Layer 3 unconditional (symmetry-only, no (H5$_F$) dependency)
and supersede the THEO-DSL-8 $\Cs$/3D identification.

\section*{Verification}
\texttt{series\_umbrella/series\_substrate\_chirality\_arc/dynamical\_substrate\_law/code/verify\_face\_alpha2\_closure.py}
builds the explicit 120-vertex 600-cell, performs the vertex cross-check
($-9/\phi^2$), verifies $\Vfour=\{e,\sigma_E,\sigma_{\!F},C_2\}$ are automorphisms
fixing $(\vhost,\nperpF)$ with $2$D invariant subspace, confirms the THEO-DSL-8 basis
degeneracy, computes the face closure $(-7/\phi^2,\,2\phi-6)$, reports the
B.1--B.4 decomposition, and confirms robustness over all $30$ incident faces. All
checks pass; coefficients exact in $\mathbb{Q}[\phi]$ with numerical agreement
$\sim 10^{-16}$.

% BEGIN GENERATED IDENTIFIER APPENDIX -- do not edit by hand
\clearpage
\section*{Appendix: Programme identifiers used in this paper}
\addcontentsline{toc}{section}{Appendix: Programme identifiers}

\noindent This programme tracks its unresolved questions under short identifiers so that a claim and its open dependencies can be cited precisely. The identifiers appearing in this paper are glossed below; no external document is needed to read them.

\begin{description}
  \item[\texttt{OPEN-FP-F1-2}] \hfill \\ Derive Mechanism A (F.1 framework axioms MA.1 + MA.2: propagation-rate asymmetry \$r(\textbackslash\{\}hat\{e\}) = r\_0(1 + \textbackslash\{\}delta\textbackslash\{\},\textbackslash\{\}hat\{e\}\textbackslash\{\}cdot\textbackslash\{\}hat\{n\})\$ and framework-local current construction at \$\textbackslash\{\}mathcal\{O\}(\textbackslash\{\}delta\textasciicircum{}1)\$) from CPP primitive axioms A1–A11 alone.
  \item[\texttt{OPEN-FP-F1-5}] \hfill \\ Extend or reformulate F.1 Theorems 5.1 (host-first-shell projection), 5.2 (first-shell perpendicularity), and 7.1 (substrate-locality umbrella) under edge-aligned Reading C (\$\textbackslash\{\}hat\{n\}\$ along a 600-cell short-edge direction; residual \$D\_5\$ symmetry at edge midpoint per Patch 0596 Remark 5.1 anti-erasure correction — see below) and face-aligned Reading C (\$\textbackslash\{\}hat\{n\}\$ perpendicular to a triangular face; residual \$C\_s\$ symmetry under vertex-host convention per Patch 0597 Remark 7.1 anti-erasure correction — see below).
\end{description}
% END GENERATED IDENTIFIER APPENDIX

\end{document}
